\documentclass[11pt,a4paper]{article}
\usepackage[utf8]{inputenc}
\usepackage[T1]{fontenc}
\usepackage[spanish]{babel}
\usepackage{geometry}
\geometry{margin=2.5cm}
\usepackage{amsmath,amssymb}
\usepackage{enumitem}
\usepackage{booktabs}
\usepackage{array}
\usepackage{longtable}
\usepackage{hyperref}
\usepackage{xcolor}
\usepackage{titlesec}
\usepackage{fancyhdr}
\usepackage{listings}
\usepackage{mdframed}
\usepackage{tcolorbox}

\definecolor{done}{RGB}{34,139,34}
\definecolor{codeblue}{RGB}{0,80,160}
\definecolor{codegray}{RGB}{240,240,240}
\definecolor{highlight}{RGB}{255,250,205}

\lstset{
  basicstyle=\ttfamily\small,
  backgroundcolor=\color{codegray},
  breaklines=true,
  frame=single,
  framesep=4pt,
  xleftmargin=6pt,
  xrightmargin=6pt,
  aboveskip=6pt,
  belowskip=6pt,
}

\pagestyle{fancy}
\fancyhf{}
\rhead{TFG --- Implementación y Pruebas}
\lhead{Mario García Blasco}
\rfoot{\thepage}

\title{\textbf{Búsqueda Semántica de Objetos en Simulación}\\[0.3cm]
\large Documento de implementación y guía de pruebas\\[0.2cm]
\normalsize Heatmaps precisos · Filtro de coherencia · Mapper LLM · YOLOE · Métodos M1--M3}
\author{Mario García Blasco}
\date{26 de marzo de 2026}

\begin{document}
\maketitle
\tableofcontents
\newpage

% ============================================================================
\section{Resumen del sistema}
% ============================================================================

El sistema implementa un pipeline de tres niveles para la búsqueda de objetos en
entornos simulados Matterport3D con Habitat-Sim:

\begin{enumerate}
  \item \textbf{Nivel 1 -- Habitaciones (LabelMe):} polígonos de habitaciones etiquetados
        manualmente que permiten al robot navegar hacia una habitación concreta.
  \item \textbf{Nivel 2 -- Mobiliario (VLMaps):} mapa semántico con 37 categorías MP3D
        que produce heatmaps de dónde es probable encontrar un mueble.
  \item \textbf{Nivel 3 -- Verificación visual (YOLOE):} detector open-vocabulary que
        confirma si el objeto específico está presente al llegar al destino.
\end{enumerate}

El fichero principal es:
\begin{lstlisting}
third_party/vlmaps/application/interactive_object_nav.py
\end{lstlisting}

% ============================================================================
\section{Características implementadas}
% ============================================================================

\subsection{Heatmaps precisos (pipeline de index\_map)}

\textbf{Problema original:} La versión de heatmap usada en navegación interactiva
re-normalizaba el mapa CLIP, produciendo regiones muy ruidosas y difusas con
falsos positivos.

\textbf{Solución:} Se reemplazó el cálculo por el mismo pipeline que usa
\texttt{index\_map.py} del repositorio original:
\begin{enumerate}
  \item \texttt{vlmap.init\_categories(mp3dcat[1:-1])} --- inicializa scores CLIP para las 37 categorías.
  \item \texttt{vlmap.index\_map(category, with\_init\_cat=True)} --- genera máscara booleana 3D.
  \item \texttt{pool\_3d\_label\_to\_2d(mask, grid\_pos, gs)} --- proyecta a vista cenital 2D.
  \item \texttt{get\_heatmap\_from\_mask\_2d(mask\_2d, cell\_size=0.05, decay\_rate=0.01)} --- genera heatmap suavizado.
\end{enumerate}

\textbf{Función clave:} \texttt{compute\_heatmap(robot, category)} en
\texttt{interactive\_object\_nav.py}.

\textbf{Efecto medido:}
\begin{center}
\begin{tabular}{lcc}
\toprule
Categoría & Celdas $>0.5$ (antiguo) & Celdas $>0.5$ (nuevo) \\
\midrule
sofa       & 22\,000 & 8\,000  \\
table      & 36\,000 & 11\,000 \\
toilet     & 18\,000 & 5\,000  \\
\bottomrule
\end{tabular}
\end{center}
El nuevo heatmap es 2--3$\times$ más concentrado en las zonas reales.

% -----------------------------------------------------------------------
\subsection{Filtro de coherencia del heatmap}

Si el heatmap está muy fragmentado (el objeto apareció en muchas celdas
dispersas), no es fiable. Se descarta y se usa fallback a BFS.

\textbf{Función:} \texttt{\_heatmap\_is\_coherent(mask\_2d, min\_total\_cells=30,
min\_dominance=0.30)}

\textbf{Criterios:}
\begin{itemize}
  \item Al menos 30 celdas activas en total.
  \item La componente conexa más grande ocupa $\geq 30\%$ del total de celdas
        activas (dominio de la componente principal).
\end{itemize}

\textbf{Resultados medidos:}
\begin{center}
\begin{tabular}{lcc}
\toprule
Categoría & Dominancia & ¿Aceptado? \\
\midrule
sofa         & 44\% & \textcolor{done}{\checkmark} \\
bed          & 41\% & \textcolor{done}{\checkmark} \\
toilet       & 45\% & \textcolor{done}{\checkmark} \\
wall         & 13\% & \textcolor{red}{$\times$} \\
chair        & 12\% & \textcolor{red}{$\times$} \\
table        & 16\% & \textcolor{red}{$\times$} \\
\bottomrule
\end{tabular}
\end{center}

% -----------------------------------------------------------------------
\subsection{Mapper LLM: traducción de objetos a categorías VLMap}

\textbf{Problema:} VLMaps solo conoce 37 categorías MP3D (muebles genéricos).
Si el usuario pide ``laptop'', VLMaps no tiene esa categoría.

\textbf{Solución:} Un LLM (GPT con few-shot) traduce el objeto a una lista de
categorías VLMap donde probablemente se encuentre.

\textbf{Función:} \texttt{map\_object\_to\_vlmap\_categories(object\_name)} en
\texttt{vlmaps/utils/llm\_utils.py}

\textbf{Ejemplos de mapeo:}
\begin{center}
\begin{tabular}{ll}
\toprule
Objeto pedido & Categorías VLMap devueltas \\
\midrule
laptop        & table, counter, shelving \\
book          & shelving, table, bed \\
mug           & counter, table \\
toilet paper  & toilet, shelving \\
tv            & tv\_monitor, furniture \\
sofa          & sofa (directo, ya está en MP3D) \\
\bottomrule
\end{tabular}
\end{center}

\textbf{Lógica:}
\begin{enumerate}
  \item Si el objeto ya es una categoría MP3D $\to$ devuelve directamente esa categoría.
  \item Si no $\to$ llama al LLM con few-shot prompting y parsea JSON
        \texttt{\{categories: [...], reason: "..."\}}.
  \item Fallback: si el parsing falla $\to$ devuelve \texttt{["furniture", "objects"]}.
\end{enumerate}

% -----------------------------------------------------------------------
\subsection{YOLOE: verificador visual en llegada}

Al llegar a un candidato, el robot gira 360° y toma fotos con la cámara
frontal. YOLOE verifica si el objeto original está visible.

\textbf{Modelo:} YOLOE-11L-Seg cargado desde \texttt{/workspace/yoloe-11l-seg.pt}

\textbf{Funciones:}
\begin{itemize}
  \item \texttt{\_load\_yoloe()} --- carga el modelo una sola vez (singleton).
  \item \texttt{yoloe\_verify(robot, object\_name, conf\_thresh=0.25)} --- captura
        frames frontales, ejecuta YOLOE con \texttt{set\_classes([object\_name])},
        devuelve \texttt{\{found, detections, frame\_bgr\}}.
\end{itemize}

\textbf{API de YOLOE (importante):}
\begin{lstlisting}[language=Python]
model.set_classes(["laptop"])   # fijar clase ANTES de predict
results = model.predict(frame, conf=0.25)
\end{lstlisting}
\textit{Nota:} la primera llamada a \texttt{set\_classes} descarga el backbone
MobileCLIP (\textasciitilde572\,MB) -- solo ocurre una vez.

\textbf{Resultados en escena real (6200 frames):}
\begin{center}
\begin{tabular}{lcc}
\toprule
Objeto & Conf. media & ¿Detectado correctamente? \\
\midrule
sofa & 0.258 & \textcolor{done}{\checkmark} \\
bed  & 0.283 & \textcolor{done}{\checkmark} \\
chair & 0.31 & \textcolor{done}{\checkmark} \\
\bottomrule
\end{tabular}
\end{center}

% -----------------------------------------------------------------------
\subsection{Análisis de regiones del heatmap (extract\_heatmap\_regions)}

Segmenta el heatmap en regiones candidatas y las ordena por calidad.

\textbf{Función:} \texttt{extract\_heatmap\_regions(heatmap, score\_thresh=0.3,
min\_area=20, min\_mean\_score=0.3, min\_density=0.4)}

\textbf{Para cada región calcula:}
\begin{itemize}
  \item \textbf{área}: número de celdas activas.
  \item \textbf{mean\_score}: media del heatmap en esas celdas.
  \item \textbf{density}: fracción de celdas activas vs.\ bounding box.
  \item \textbf{quality}: $0.3 \cdot \hat{a} + 0.4 \cdot \bar{s} + 0.3 \cdot d$
        donde $\hat{a}$ es área normalizada.
\end{itemize}

\textbf{Devuelve:} lista de dicts ordenada por \texttt{quality} descendente:
\begin{lstlisting}[language=Python]
[{"centroid": (r, c), "area": 120, "mean_score": 0.62,
  "density": 0.78, "quality": 0.71}, ...]
\end{lstlisting}

\textbf{Resultados medidos:}
\begin{center}
\begin{tabular}{lcc}
\toprule
Categoría & Nº regiones & Mejor quality \\
\midrule
sofa       & 27 & 0.809 \\
toilet     &  6 & 0.824 \\
tv\_monitor & 16 & 0.861 \\
\bottomrule
\end{tabular}
\end{center}

% -----------------------------------------------------------------------
\subsection{Métodos de búsqueda M1, M2, M3}

Al iniciar la navegación interactiva, el sistema pregunta qué método usar.

\begin{center}
\begin{tabular}{clp{7cm}}
\toprule
Método & Nombre & Descripción \\
\midrule
M1 & Random      & Selecciona un punto navegable aleatorio. Baseline ingenuo. \\
M2 & Nearest     & BFS: candidato más cercano al robot según el heatmap. Si el heatmap no es coherente, usa fallback aleatorio. \\
M3 & Semantic    & Usa \texttt{extract\_heatmap\_regions} y va directamente a la región de mayor calidad. \\
\bottomrule
\end{tabular}
\end{center}

\textbf{Funciones:}
\begin{itemize}
  \item \texttt{select\_goal\_random(robot)} --- M1.
  \item \texttt{select\_goal\_nearest(robot, furniture\_cat, heatmap)} --- M2.
  \item \texttt{select\_goal\_semantic(robot, furniture\_cat, heatmap)} --- M3.
  \item \texttt{\_METHOD\_SELECTORS = \{"M1": ..., "M2": ..., "M3": ...\}} --- dispatch.
  \item \texttt{\_active\_method = "M2"} --- valor por defecto.
\end{itemize}

% ============================================================================
\section{Cómo probar cada característica}
% ============================================================================

\subsection{Requisitos previos}

\begin{enumerate}
  \item Iniciar el contenedor Docker con GPU:
\begin{lstlisting}
./docker/run_docker.sh
\end{lstlisting}
  \item Desde el menú principal (\texttt{./docker/menu.sh}):
        \begin{itemize}
          \item Opción \textbf{b)} para construir el mapa VLMap (si no existe).
          \item Opción \textbf{n)} para iniciar navegación interactiva.
        \end{itemize}
\end{enumerate}

% -----------------------------------------------------------------------
\subsection{Probar los heatmaps precisos}

Al pedir un objeto en navegación interactiva, se imprime en consola:

\begin{lstlisting}
[Heatmap] sofa: 8234 cells > 0
[Heatmap] coherence: dominant_ratio=0.44 -> ACCEPTED
\end{lstlisting}

Para comparar con el método antiguo, comentar temporalmente el bloque de
\texttt{vlmap.init\_categories} en \texttt{compute\_heatmap()} y ver que el
número de celdas activas se multiplica por 2-3.

% -----------------------------------------------------------------------
\subsection{Probar el filtro de coherencia}

Pedir al robot que busque ``wall'' o ``floor'':

\begin{lstlisting}
You: find the wall
[Heatmap] wall: 45123 cells > 0
[Heatmap] coherence: dominant_ratio=0.13 -> REJECTED (fallback)
[Nav] Using random fallback goal
\end{lstlisting}

Contraste: pedir ``sofa'' o ``toilet'' -- el heatmap es aceptado y el robot va directamente.

% -----------------------------------------------------------------------
\subsection{Probar el Mapper LLM}

En la sesión interactiva, pedir un objeto no-MP3D:

\begin{lstlisting}
You: find the laptop
[Mapper LLM] laptop -> ['table', 'counter', 'shelving'] (reason: laptops are typically on tables or counters)
[Nav] Searching for furniture categories: table, counter, shelving
\end{lstlisting}

Pedir un objeto MP3D directo para ver el path rápido:

\begin{lstlisting}
You: find the sofa
[Mapper LLM] sofa is already an MP3D category -> ['sofa']
\end{lstlisting}

% -----------------------------------------------------------------------
\subsection{Probar YOLOE}

Al llegar al primer candidato, YOLOE ejecuta la verificación:

\begin{lstlisting}
[YOLOE] Verifying 'laptop' at arrival...
[YOLOE] FOUND: laptop (conf=0.41) -> SUCCESS
\end{lstlisting}

O si no lo encuentra:

\begin{lstlisting}
[YOLOE] laptop NOT found at this location, continuing search...
\end{lstlisting}

La ventana de OpenCV muestra el frame con bounding boxes cuando hay detecciones.

\textbf{Test unitario rápido} (sin Habitat, solo verificar que YOLOE carga):

\begin{lstlisting}[language=Python]
from ultralytics import YOLOE
model = YOLOE("/workspace/yoloe-11l-seg.pt")
model.set_classes(["sofa"])
import numpy as np
img = np.zeros((480, 640, 3), dtype=np.uint8)
results = model.predict(img, conf=0.25, verbose=False)
print("YOLOE OK, detections:", len(results[0].boxes))
\end{lstlisting}

% -----------------------------------------------------------------------
\subsection{Probar el análisis de regiones}

En el REPL de Python dentro del contenedor:

\begin{lstlisting}[language=Python]
import sys
sys.path.insert(0, "/workspace/vlmaps")
from application.interactive_object_nav import (
    compute_heatmap, extract_heatmap_regions
)
# Asumiendo que robot ya está inicializado en una sesión de navegación,
# o cargando un heatmap guardado:
regions = extract_heatmap_regions(heatmap, score_thresh=0.3, min_area=20)
for r in regions[:3]:
    print(f"  centroid={r['centroid']}  quality={r['quality']:.3f}")
\end{lstlisting}

% -----------------------------------------------------------------------
\subsection{Probar los métodos M1, M2, M3}

Al iniciar la sesión interactiva (\texttt{menu.sh} opción n):

\begin{lstlisting}
Search methods: M1 (random), M2 (nearest), M3 (semantic quality)
Select method [M1/M2/M3, default M2]: M3
\end{lstlisting}

\textbf{Prueba comparativa sugerida:}
\begin{enumerate}
  \item Escena 2, objeto: ``sofa''.
  \item Ejecutar 3 veces con M1, anotar pasos hasta encontrarlo.
  \item Ejecutar 3 veces con M2, comparar.
  \item Ejecutar 3 veces con M3, comparar.
  \item M3 debería requerir menos pasos que M1, con M2 como intermedio.
\end{enumerate}

% ============================================================================
\section{Pipeline completo: ejemplo de sesión}
% ============================================================================

\begin{lstlisting}
$ ./docker/menu.sh
> n) Interactive navigation

Search methods: M1 (random), M2 (nearest), M3 (semantic quality)
Select method [M1/M2/M3, default M2]: M3

Scene 2 loaded. Robot at (362, 480).

[M3] You: find the laptop

[Mapper LLM] laptop -> ['table', 'counter', 'shelving']

[Heatmap] table: 11230 cells > 0
[Heatmap] coherence: dominant_ratio=0.38 -> ACCEPTED
[Regions] table: 12 regions found, best quality=0.71

[Nav] Going to table region at (401, 512) [quality=0.71]
... (robot navega) ...

[YOLOE] Verifying 'laptop' at arrival...
[YOLOE] laptop NOT found, trying next region...

[Nav] Going to table region at (388, 495) [quality=0.68]
... (robot navega) ...

[YOLOE] FOUND: laptop (conf=0.37) -> SUCCESS
Goal reached in 2 attempts, 184 steps.
\end{lstlisting}

% ============================================================================
\section{Estructura de ficheros relevantes}
% ============================================================================

\begin{center}
\begin{longtable}{p{7.5cm}p{7cm}}
\toprule
\textbf{Fichero} & \textbf{Cambios} \\
\midrule
\texttt{third\_party/vlmaps/application/} \texttt{interactive\_object\_nav.py}
& Todas las funciones nuevas: \texttt{compute\_heatmap}, \texttt{\_heatmap\_is\_coherent}, \texttt{find\_heatmap\_goal}, \texttt{extract\_heatmap\_regions}, \texttt{\_load\_yoloe}, \texttt{yoloe\_verify}, \texttt{select\_goal\_random/nearest/semantic}, \texttt{\_METHOD\_SELECTORS} \\
\midrule
\texttt{third\_party/vlmaps/vlmaps/utils/} \texttt{llm\_utils.py}
& \texttt{map\_object\_to\_vlmap\_categories()} con few-shot LLM \\
\midrule
\texttt{third\_party/vlmaps/vlmaps/robot/} \texttt{habitat\_lang\_robot.py}
& \texttt{\_NAV\_MARGIN\_PX = 2} (margen de 2px en mapa erosionado $\approx$10\,cm) \\
\bottomrule
\end{longtable}
\end{center}

% ============================================================================
\section{Próximos pasos}
% ============================================================================

\begin{enumerate}
  \item \textbf{M4}: scoring híbrido (VLMap $\times$ YOLOE confidence) con
        re-routing si YOLOE no detecta tras $k$ visitas.
  \item \textbf{M5}: planificación de 4 viewpoints por región + exploración de
        fronteras para regiones no alcanzadas.
  \item \textbf{Benchmark}: \texttt{run\_evaluation.py} con 10 objetos $\times$
        3 escenas $\times$ 5 métodos (M1--M5), métricas: success rate,
        avg steps, avg path length.
  \item \textbf{ROS 1 Noetic}: wrapper para Toyota HSR (última prioridad).
\end{enumerate}

% ============================================================================
\section{Actualización — 2 de abril de 2026}
% ============================================================================

\subsection{Postproceso de heatmaps: \texttt{postprocess\_heatmap}}

\subsubsection*{Motivación}
Al generar mapas de similitud semántica para queries como \textit{bed},
\textit{window} o \textit{bathroom}, aparece una región principal coherente
junto con múltiples activaciones pequeñas dispersas (ruido). El postproceso
las elimina de forma genérica, sin reglas por clase.

\subsubsection*{Pipeline implementado}

\begin{enumerate}
  \item \textbf{Suavizado Gaussiano} (\texttt{blur\_ksize=5},
        \texttt{blur\_sigma=1.0}): reduce ruido de alta frecuencia.
  \item \textbf{Umbral relativo}: \texttt{thr = rel\_thresh * max(smoothed)},
        con \texttt{rel\_thresh=0.5}. Adaptativo a cada query y escena.
  \item \textbf{Componentes conectadas} (8-conectividad, OpenCV
        \texttt{connectedComponentsWithStats}).
  \item \textbf{Puntuación por componente}:
        $\text{score} = \bar{v} \cdot \ln(1+A)$
        (\texttt{mean\_log\_area}), donde $\bar{v}$ es el valor medio y $A$
        el área. Alternativa: \texttt{max\_sqrt\_area}.
  \item \textbf{Filtro relativo}: se retienen componentes con
        $\text{score} \geq \texttt{keep\_ratio} \cdot \text{best\_score}$,
        con \texttt{keep\_ratio=0.2}. Preserva objetos pequeños reales con
        alta similitud.
\end{enumerate}

\textbf{API}:
\begin{verbatim}
cleaned_heatmap, final_mask, kept_components = postprocess_heatmap(
    heatmap, blur_ksize=5, blur_sigma=1.0, rel_thresh=0.5,
    min_area=3, score_mode="mean_log_area", keep_ratio=0.2)
\end{verbatim}

\texttt{compute\_heatmap} ahora devuelve
\texttt{(heatmap, kept\_components)}.

\subsubsection*{Coste computacional}
Gaussian blur + connected components sobre un mapa $1000\times1000$ en
NumPy/OpenCV: $\approx$5--15~ms en CPU. Apto para uso interactivo.

% -----------------------------------------------------------------------
\subsection{Exploración 360° tras fallo de detección}
% -----------------------------------------------------------------------

\subsubsection*{Flujo}
Después de Stage~1 (0{,}5~m) y Stage~2 (1~m) sin confirmación YOLOE, se
ejecuta \texttt{scan\_360\_and\_verify}:

\begin{enumerate}
  \item El robot rota $360°$ en 18 pasos de 20° cada uno.
  \item Cada paso usa \texttt{robot.turn(20)}, que internamente ejecuta
        cuatro acciones discretas \texttt{turn\_right} de 5°.
        \textbf{No hay teletransporte ni saltos bruscos.}
  \item Tras cada giro se captura el frame de cámara y se ejecuta YOLOE.
  \item Si YOLOE confirma el objeto, la exploración termina.
\end{enumerate}

% -----------------------------------------------------------------------
\subsection{Ruta alternativa (\texttt{navigate\_to\_alternative})}
% -----------------------------------------------------------------------

Si la exploración 360° también falla, el sistema selecciona la componente
del heatmap postprocesado con mayor score cuyo centroide esté a
$\geq20$ celdas (${\approx}1$~m) de la posición actual, planifica un standoff de
1~m sobre la región del VLMap más cercana y navega hasta allí mediante
\texttt{robot.move\_to()} (acciones discretas). Tras llegar, se ejecuta
una última verificación YOLOE.

\subsubsection*{Resumen del flujo completo (Stage 1--4)}

\begin{center}
\begin{tabular}{clll}
\toprule
Stage & Posición & YOLOE & Si falla $\to$ \\
\midrule
1 & 0{,}5~m standoff & Sí & Stage 2 \\
2 & 1~m standoff     & Sí & Stage 3 \\
3 & 360° scan (misma pos.) & Sí ($\times$18) & Stage 4 \\
4 & Alt.\ componente heatmap & Sí & Desistir \\
\bottomrule
\end{tabular}
\end{center}

\subsection{Archivos modificados}

\begin{center}
\begin{tabular}{ll}
\toprule
\textbf{Archivo} & \textbf{Cambio} \\
\midrule
\texttt{application/interactive\_object\_nav.py}
  & \texttt{postprocess\_heatmap}, \texttt{scan\_360\_and\_verify},
    \texttt{navigate\_to\_alternative}; loop principal extendido con
    Stage~3 y Stage~4 \\
\bottomrule
\end{tabular}
\end{center}

\end{document}
